\chapter{Introduction to Arduino and Sensors}
\label{chap:arduino-sensors}

\section{Microcontroller Architecture and Embedded Systems}
An \textbf{embedded system} is a specialized computing system composed of a microcontroller, dedicated firmware, sensory inputs, and actuation outputs designed to execute deterministic real-time control functions.

\subsection{Architectural Classifications: Von Neumann vs Harvard}
\begin{itemize}
    \item \textbf{Von Neumann Architecture:} Employs a unified memory space and shared system bus for both program instructions and runtime data. Instruction fetch and data operand read/write operations must be executed sequentially in separate machine cycles, resulting in the \textbf{Von Neumann bus bottleneck}.
    \item \textbf{Harvard Architecture:} Employs physically separate memory spaces and independent instruction and data buses. This allows simultaneous instruction fetching and data reading/writing, significantly enhancing pipelining efficiency and instruction throughput.
\end{itemize}

\subsection{The ATmega328P Microcontroller Core}
The standard Arduino Uno development platform is built around the Microchip/Atmel **ATmega328P**, an 8-bit AVR enhanced RISC microcontroller.

\begin{table}[htbp]
\centering
\caption{Key Hardware Specifications of the ATmega328P}
\label{tab:atmega328p-specs}
\begin{tabularx}{\textwidth}{l X}
\toprule
\textbf{Parameter / Subsystem} & \textbf{Specification / Characteristics} \\
\midrule
CPU Architecture & 8-bit AVR RISC (131 instructions, most single clock cycle execution) \\
Operating Frequency ($f_{osc}$) & $\SI{16}{\mega\hertz}$ (provides execution throughput of 16 MIPS) \\
Operating Voltage ($V_{cc}$) & $\SI{1.8}{\volt}$ to $\SI{5.5}{\volt}$ (configured at $\SI{5.0}{\volt}$ on Arduino Uno) \\
Flash Program Memory & $\SI{32}{\kibi\byte}$ In-System Programmable Flash ($\SI{0.5}{\kibi\byte}$ bootloader) \\
SRAM (Static RAM) & $\SI{2}{\kibi\byte}$ internal volatile data memory \\
EEPROM & $\SI{1}{\kibi\byte}$ non-volatile byte-addressable parameter memory \\
General Purpose Registers & 32 $\times$ 8-bit general-purpose working registers ($R_0 \dots R_{31}$) \\
GPIO Ports & Port B (PB0--PB5), Port C (PC0--PC5), Port D (PD0--PD7) \\
Analog ADC Channels & 6 multiplexed 10-bit Successive Approximation (SAR) channels \\
Timer/Counter Modules & Two 8-bit timers (Timer0, Timer2) and one 16-bit timer (Timer1) \\
Hardware PWM Channels & 6 channels ($\text{D3}, \text{D5}, \text{D6}, \text{D9}, \text{D10}, \text{D11}$) \\
\bottomrule
\end{tabularx}
\end{table}

\subsection{Low-Level GPIO Register Architecture}
Each physical 8-bit I/O port (Port B, Port C, Port D) is controlled by three dedicated memory-mapped control registers:
\begin{enumerate}
    \item \textbf{Data Direction Register (\texttt{DDRx}):} Configures pin directions ($0 = \text{Input}, 1 = \text{Output}$).
    \item \textbf{Port Output Data Register (\texttt{PORTx}):} Sets output voltage level when configured as output ($0 = \text{LOW}, 1 = \text{HIGH}$) or enables internal $\SI{20}{\kilo\ohm}\text{--}\SI{50}{\kilo\ohm}$ pull-up resistors when configured as input.
    \item \textbf{Port Input Pins Address (\texttt{PINx}):} Read-only register reflecting the immediate physical electrical logic levels present at the port pins.
\end{enumerate}

\section{Analog-to-Digital Conversion (ADC)}
The ATmega328P incorporates an internal 10-bit Successive Approximation Register (SAR) ADC.

\subsection{Mathematical Principles and Quantization}
A 10-bit ADC quantizes a continuous analog input voltage $V_{in}$ into one of $2^{10} = 1024$ discrete digital levels ($0 \dots 1023$).

\begin{derivationbox}[ADC Resolution and Voltage Calculation]
The voltage resolution (Least Significant Bit step size $\Delta V$) is defined by:
\begin{equation}
\Delta V = \frac{V_{ref}}{2^n} = \frac{V_{ref}}{1024}
\end{equation}
When using the default internal $\SI{5.0}{\volt}$ supply reference:
\begin{equation}
\Delta V = \frac{\SI{5.0}{\volt}}{1024} = \SI{4.8828}{\milli\volt}/\text{LSB}
\end{equation}
The digital integer output code $D_{out} \in [0, 1023]$ is given by:
\begin{equation}
D_{out} = \left\lfloor \frac{V_{in}}{V_{ref}} \times 1024 \right\rfloor \implies V_{in} = \frac{D_{out} \times V_{ref}}{1023}
\end{equation}
\end{derivationbox}

The ADC conversion circuitry requires an input clock between $\SI{50}{\kilo\hertz}$ and $\SI{200}{\kilo\hertz}$. For a $\SI{16}{\mega\hertz}$ system clock, an internal prescaler of 128 yields an ADC clock of $\SI{125}{\kilo\hertz}$, resulting in a total conversion time of 13 ADC clock cycles ($\approx \SI{104}{\micro\second}$ per sample).

\section{Pulse Width Modulation (PWM)}
Pulse Width Modulation is a technique for generating pseudo-analog effective output voltages from digital switching pins by varying the duty cycle of a fixed-frequency rectangular pulse train.

\subsection{Mathematical Foundations of PWM}
\begin{itemize}
    \item \textbf{Duty Cycle ($D$):} The ratio of active-high on-time ($t_{on}$) to total switching period ($T$):
    \begin{equation}
    D = \frac{t_{on}}{T} \times 100\% = \frac{t_{on}}{t_{on} + t_{off}} \times 100\%
    \end{equation}
    \item \textbf{Average DC Voltage ($V_{avg}$):}
    \begin{equation}
    V_{avg} = \frac{1}{T} \int_0^T v(t)\,dt = D \times V_{cc} = \left( \frac{\text{OCR}}{255} \right) \times V_{cc}
    \end{equation}
    where $\text{OCR}$ is the 8-bit Output Compare Register value ($0 \le \text{OCR} \le 255$).
\end{itemize}

On Arduino Uno, pins 3, 5, 6, 9, 10, and 11 support hardware PWM. Pins 5 and 6 operate at a base frequency of $\SI{980}{\hertz}$ (driven by Timer0), while pins 3, 9, 10, and 11 operate at $\SI{490}{\hertz}$ (driven by Timer1 and Timer2).

\section{Hardware Interrupts and Real-Time Event Handling}
\subsection{Polling vs Interrupts}
\begin{itemize}
    \item \textbf{Polling:} The CPU continuously executes software loops to test I/O pin states. Polling wastes processor cycles and introduces variable latency depending on loop execution depth.
    \item \textbf{Hardware Interrupts:} An asynchronous electrical edge at an external pin immediately suspends normal CPU execution, preserves program counter context on the stack, and branches directly to a dedicated \textbf{Interrupt Service Routine (ISR)}.
\end{itemize}

\subsection{Interrupt Architecture on ATmega328P}
Arduino Uno provides two external hardware interrupt pins:
\begin{itemize}
    \item \textbf{INT0} mapped to Digital Pin 2.
    \item \textbf{INT1} mapped to Digital Pin 3.
\end{itemize}
Interrupt triggers can be configured for four distinct hardware modes: \texttt{LOW} (pin is low), \texttt{CHANGE} (any logic transition), \texttt{RISING} ($0 \to 1$ edge), or \texttt{FALLING} ($1 \to 0$ edge).

\begin{pitfallbox}[ISR Design Rules and the \texttt{volatile} Keyword]
\begin{enumerate}
    \item ISR execution must be kept extremely brief (e.g., updating a flag or incrementing a counter).
    \item Do \textbf{not} invoke blocking delays (\texttt{delay()}) or complex serial transmissions (\texttt{Serial.print()}) inside an ISR, as hardware interrupts are globally disabled during ISR execution.
    \item Any global variable modified inside an ISR and read in the main loop \textbf{must} be qualified with the \texttt{volatile} keyword to instruct the compiler optimizer not to cache the variable in CPU registers:
    \begin{lstlisting}[language=C++]
    volatile unsigned long pulseCount = 0;
    \end{lstlisting}
\end{enumerate}
\end{pitfallbox}

\section{Sensors, Transducers, and Actuators}
\subsection{Light Dependent Resistor (LDR)}
An LDR (cadmium sulfide photoresistor) exhibits photoconductivity: photon absorption excites valence electrons into the conduction band, decreasing electrical resistance ($R_{LDR} \approx \SI{1}{\mega\ohm}$ in total darkness down to $R_{LDR} \approx \SI{500}{\ohm}$ in bright sunlight).

\begin{derivationbox}[LDR Voltage Divider Analysis]
Connecting the LDR in series with a fixed pull-down resistor $R_{fixed} = \SI{10}{\kilo\ohm}$ forms a voltage divider:
\begin{equation}
V_{out} = V_{cc} \left( \frac{R_{fixed}}{R_{fixed} + R_{LDR}} \right)
\end{equation}
As ambient illumination increases, $R_{LDR}$ decreases, causing $V_{out}$ to rise monotonically towards $V_{cc}$.
\end{derivationbox}

\subsection{Infrared (IR) Obstacle Sensor}
An active reflective IR sensor comprises an infrared emitting diode (wavelength $\lambda \approx \SI{940}{\nano\second}$) and an IR photodiode receiver coupled to an on-board LM393 voltage comparator.
When an obstacle enters the detection field, reflected IR energy increases reverse photodiode current, causing the non-inverting comparator terminal to cross the reference threshold set by a precision potentiometer, driving the digital output (\texttt{OUT}) active \textbf{LOW}.

\subsection{Ultrasonic Distance Sensor (HC-SR04)}
The HC-SR04 measures physical distance using non-contact acoustic echo-ranging.

\begin{derivationbox}[Ultrasonic Range Derivation]
The measurement sequence:
\begin{enumerate}
    \item A $\SI{10}{\micro\second}$ active-HIGH pulse is applied to the \texttt{Trig} pin.
    \item The module generates an 8-cycle ultrasonic acoustic burst at $\SI{40}{\kilo\hertz}$.
    \item The \texttt{Echo} pin goes HIGH and remains HIGH for the duration $t_{echo}$ taken by the acoustic wave to travel to the target and reflect back.
    \item Given the acoustic velocity in air at $\SI{20}{\celsius}$ is $v_{sound} \approx \SI{343}{\meter\per\second} = \SI{0.0343}{\centi\meter\per\micro\second}$:
    \begin{equation}
    d = \frac{v_{sound} \times t_{echo}}{2} = \frac{t_{echo} \text{ (\SI{}{\micro\second})}}{58.3} \text{ \SI{}{\centi\meter}}
    \end{equation}
\end{enumerate}
\end{derivationbox}

\subsection{Digital Temperature and Humidity Sensor (DHT11/DHT22)}
The DHT11 integrates a capacitive relative humidity sensor, an NTC thermistor, and an internal 8-bit microcontroller that outputs data over a single-wire bi-directional serial bus.
\begin{itemize}
    \item \textbf{Data Framing (40 Bits):} Transmits 16 bits integral/decimal relative humidity, 16 bits integral/decimal temperature, and an 8-bit checksum:
    \begin{equation}
    \text{Checksum} = (\text{Byte}_1 + \text{Byte}_2 + \text{Byte}_3 + \text{Byte}_4) \ \& \ \text{0xFF}
    \end{equation}
\end{itemize}

\subsection{Motor Actuators and Driver Interfaces}
\begin{itemize}
    \item \textbf{DC Motor Control via L293D H-Bridge:} Direct connection of DC motors to microcontroller GPIO pins will destroy the chip due to excessive current draw and inductive flyback voltages. The L293D quadruple high-current H-bridge IC provides bidirectional drive up to $\SI{600}{\milli\ampere}$ per channel with internal clamp diodes.
    \item \textbf{Servo Motor Control:} Standard hobby servomotors employ PWM pulse-width positioning operating at a frame frequency of $\SI{50}{\hertz}$ ($T = \SI{20}{\milli\second}$):
    \begin{itemize}
        \item $\SI{1.0}{\milli\second}$ high pulse $\implies 0^\circ$ shaft position.
        \item $\SI{1.5}{\milli\second}$ high pulse $\implies 90^\circ$ neutral position.
        \item $\SI{2.0}{\milli\second}$ high pulse $\implies 180^\circ$ maximum position.
    \end{itemize}
\end{itemize}

\section{Complete Production Arduino Firmware Implementations}

\subsection{Multi-Sensor Acquisition and Display Firmware}
\begin{lstlisting}[language=C++]
// Complete Production-Grade Arduino Multi-Sensor Firmware
#include <Arduino.h>

const uint8_t PIN_LDR = A0;
const uint8_t PIN_TRIG = 9;
const uint8_t PIN_ECHO = 10;
const uint8_t PIN_LED_PWM = 6;
const uint8_t PIN_INTERRUPT = 2;

volatile unsigned long tachPulseCount = 0;
const float V_REF = 5.0;

void isrTachometer() {
    tachPulseCount++;
}

float measureDistanceCm() {
    digitalWrite(PIN_TRIG, LOW);
    delayMicroseconds(2);
    digitalWrite(PIN_TRIG, HIGH);
    delayMicroseconds(10);
    digitalWrite(PIN_TRIG, LOW);
    
    unsigned long duration = pulseIn(PIN_ECHO, HIGH, 30000); // 30ms timeout
    if (duration == 0) return -1.0; // Out of range
    return (duration * 0.0343) / 2.0;
}

void setup() {
    Serial.begin(115200);
    pinMode(PIN_TRIG, OUTPUT);
    pinMode(PIN_ECHO, INPUT);
    pinMode(PIN_LED_PWM, OUTPUT);
    pinMode(PIN_INTERRUPT, INPUT_PULLUP);
    
    attachInterrupt(digitalPinToInterrupt(PIN_INTERRUPT), isrTachometer, FALLING);
    Serial.println(F("System Initialized Successfully."));
}

void loop() {
    // 1. Read Analog Light Level
    int rawADC = analogRead(PIN_LDR);
    float vSensor = (rawADC * V_REF) / 1023.0;
    
    // 2. Adjust PWM brightness proportionally (inverted: darker = brighter LED)
    uint8_t pwmValue = map(rawADC, 0, 1023, 255, 0);
    analogWrite(PIN_LED_PWM, pwmValue);
    
    // 3. Ultrasonic Distance
    float dist = measureDistanceCm();
    
    // 4. Safely read volatile interrupt counter
    noInterrupts();
    unsigned long currentPulses = tachPulseCount;
    interrupts();
    
    // 5. Formatted Telemetry Output
    Serial.print(F("ADC: ")); Serial.print(rawADC);
    Serial.print(F(" | V: ")); Serial.print(vSensor, 2);
    Serial.print(F("V | Dist: ")); Serial.print(dist, 1);
    Serial.print(F("cm | Pulses: ")); Serial.println(currentPulses);
    
    delay(500);
}
\end{lstlisting}

\section{Worked Numerical Examples}

\begin{examplebox}[Example 6.1: Precision ADC Voltage Calculation]
\textbf{Problem:} An ATmega328P ADC powered at $V_{ref} = \SI{5.00}{\volt}$ returns an integer readout of 682 from an analog sensor. Determine:
\begin{enumerate}
    \item The exact measured analog input voltage $V_{in}$.
    \item The quantization error bound.
\end{enumerate}

\textbf{Solution:}
\begin{enumerate}
    \item \textbf{Input Voltage Calculation:}
    \begin{equation}
    V_{in} = \frac{D_{out} \times V_{ref}}{1023} = \frac{682 \times \SI{5.00}{\volt}}{1023} \approx \SI{3.333}{\volt}
    \end{equation}
    \item \textbf{Quantization Error:}
    The step size $\Delta V = \frac{\SI{5.00}{\volt}}{1024} = \SI{4.883}{\milli\volt}$.
    The maximum quantization uncertainty is $\pm \frac{1}{2} \Delta V = \pm \SI{2.441}{\milli\volt}$.
\end{enumerate}
\end{examplebox}

\begin{examplebox}[Example 6.2: PWM Duty Cycle and Equivalent DC Level]
\textbf{Problem:} An 8-bit Arduino PWM timer is configured with an Output Compare value of $\text{OCR} = 178$ operating from a $\SI{5.0}{\volt}$ logic supply. Calculate the duty cycle percentage $D$ and the equivalent average DC output voltage $V_{avg}$.

\textbf{Solution:}
\begin{enumerate}
    \item \textbf{Duty Cycle Calculation:}
    \begin{equation}
    D = \left( \frac{\text{OCR}}{255} \right) \times 100\% = \left( \frac{178}{255} \right) \times 100\% \approx 69.80\%
    \end{equation}
    \item \textbf{Average DC Voltage Calculation:}
    \begin{equation}
    V_{avg} = D \times V_{cc} = 0.6980 \times \SI{5.0}{\volt} = \SI{3.490}{\volt}
    \end{equation}
\end{enumerate}
\end{examplebox}

\begin{examplebox}[Example 6.3: HC-SR04 Ultrasonic Timing and Distance]
\textbf{Problem:} An HC-SR04 ultrasonic sensor measures an echo pulse width of $t_{echo} = \SI{1458}{\micro\second}$. If the ambient temperature is $\SI{20}{\celsius}$ ($v_{sound} = \SI{343}{\meter\per\second}$), determine the physical distance to the object in centimeters.

\textbf{Solution:}
\begin{enumerate}
    \item Acoustic speed in air: $v_{sound} = \SI{343}{\meter\per\second} = \SI{0.0343}{\centi\meter\per\micro\second}$.
    \item Two-way travel distance $2d = v_{sound} \times t_{echo}$.
    \item One-way distance to target:
    \begin{equation}
    d = \frac{0.0343\,\text{cm}/\mu\text{s} \times \SI{1458}{\micro\second}}{2} = \frac{\SI{50.009}{\centi\meter}}{2} \approx \SI{25.00}{\centi\meter}
    \end{equation}
\end{enumerate}
\end{examplebox}

\begin{examplebox}[Example 6.4: LDR Voltage Divider Analysis]
\textbf{Problem:} An LDR is connected in a voltage divider with a fixed resistor $R_{fixed} = \SI{10}{\kilo\ohm}$ tied to ground, powered from $V_{cc} = \SI{5.0}{\volt}$. Under ambient room lighting, the LDR exhibits a resistance of $R_{LDR} = \SI{2.5}{\kilo\ohm}$. Under direct sunlight, $R_{LDR}$ drops to $\SI{400}{\ohm}$. Calculate the voltage divider output $V_{out}$ and the corresponding 10-bit ADC integer readings for both conditions.

\textbf{Solution:}
\begin{enumerate}
    \item \textbf{Room Light Condition ($R_{LDR} = \SI{2.5}{\kilo\ohm}$):}
    \begin{align*}
    V_{out, room} &= V_{cc} \left( \frac{R_{fixed}}{R_{fixed} + R_{LDR}} \right) = \SI{5.0}{\volt} \times \left( \frac{\SI{10}{\kilo\ohm}}{\SI{10}{\kilo\ohm} + \SI{2.5}{\kilo\ohm}} \right) = \SI{5.0}{\volt} \times 0.80 = \SI{4.00}{\volt} \\
    \text{ADC reading} &= \left\lfloor \frac{4.00}{5.00} \times 1023 \right\rfloor = 818
    \end{align*}
    \item \textbf{Direct Sunlight Condition ($R_{LDR} = \SI{400}{\ohm} = \SI{0.4}{\kilo\ohm}$):}
    \begin{align*}
    V_{out, sun} &= \SI{5.0}{\volt} \times \left( \frac{\SI{10}{\kilo\ohm}}{\SI{10}{\kilo\ohm} + \SI{0.4}{\kilo\ohm}} \right) = \SI{5.0}{\volt} \times \frac{10}{10.4} \approx \SI{4.808}{\volt} \\
    \text{ADC reading} &= \left\lfloor \frac{4.808}{5.00} \times 1023 \right\rfloor = 983
    \end{align*}
\end{enumerate}
\end{examplebox}

\begin{examplebox}[Example 6.5: DHT11 Checksum Verification]
\textbf{Problem:} A DHT11 sensor transmits the following 5 raw data bytes:
\begin{itemize}
    \item $\text{Byte}_1$ (Integral RH): $00110101_2$ ($53_{10}$)
    \item $\text{Byte}_2$ (Decimal RH): $00000000_2$ ($0_{10}$)
    \item $\text{Byte}_3$ (Integral Temp): $00011000_2$ ($24_{10}$)
    \item $\text{Byte}_4$ (Decimal Temp): $00000000_2$ ($0_{10}$)
    \item $\text{Byte}_5$ (Checksum received): $01001101_2$ ($77_{10}$)
\end{itemize}
Validate if the transmission is error-free and state the measured temperature and humidity.

\textbf{Solution:}
\begin{enumerate}
    \item Calculate expected checksum:
    \begin{equation}
    \text{Sum} = 53 + 0 + 24 + 0 = 77_{10} = 01001101_2
    \end{equation}
    \item Compare with received $\text{Byte}_5$:
    $\text{Calculated Checksum } (77) == \text{Received Checksum } (77)$.
    The transmission is valid and error-free.
    \item Physical readings: Relative Humidity $= 53\%$, Temperature $= \SI{24}{\celsius}$.
\end{enumerate}
\end{examplebox}

\begin{takeawaybox}[Chapter 6 Key Formulae \& Takeaways]
\begin{itemize}
    \item \textbf{10-bit ADC Formula:} $V_{in} = \frac{\text{ADC\_val} \times V_{ref}}{1023}$.
    \item \textbf{PWM Average Voltage:} $V_{avg} = \left( \frac{\text{OCR}}{255} \right) V_{cc}$.
    \item \textbf{Ultrasonic Distance:} $d = \frac{t_{echo} \times 0.0343}{2}$ (\SI{}{\centi\meter}).
    \item \textbf{Interrupt Best Practices:} Keep ISRs non-blocking, qualify shared variables as \texttt{volatile}.
\end{itemize}
\end{takeawaybox}
